% ---------------------------------------------------------------------------
% Discussion additions
%
% Two items from the review:
%   (1) "clinically sound / bridge the gap" -> what a reader who builds clinical
%       models should actually do differently. Insert AFTER the first Discussion
%       paragraph (the one that states the findings) and BEFORE the limitations.
%   (2) "gestire dati multimodali in oncologia usando FM 2+" -> future work.
%       Insert AFTER the limitations paragraph, before the Conclusion.
%
% CITATIONS: keys marked <TODO-CITE> are placeholders. Do not compile them as
% real \cite commands until the references are verified -- the bibliography
% already carries six entries flagged [Verify].
% ---------------------------------------------------------------------------


% ===== Insert after the findings paragraph of the Discussion =================

These results carry a practical consequence for how a transfer is approached in
a multi-center clinical project. The two shift mechanisms that are measurable
without target labels --- a change in the covariate distribution and a change in
event prevalence --- are also the two that behave most differently: the second is
correctable post hoc, while the first is not the reason discrimination is lost.
A pragmatic protocol follows. Before deploying a model at a new center, measure
both quantities, since both are computable from unlabeled data. If prevalence
has moved, recalibrate, because that is the one regime in which an unsupervised
correction demonstrably restores the probabilities that a clinical decision
actually depends on. If discrimination has dropped, our results indicate that no
amount of unsupervised alignment will return it, and effort is better spent on
the causes upstream of the model: harmonizing how the predictors are measured
across centers, or acquiring a modest number of labels at the target site. The
contribution we intend here is diagnostic rather than prescriptive. A reported
transfer failure conflates two things --- a transfer that is intrinsically
unsolvable from unlabeled data, and a method that underperforms on a solvable
one --- and separating them is what tells a clinical team whether to change the
method or to change the data. The oracle makes that separation measurable on
synthetic data, and the two measurable mechanisms approximate it on real
cohorts.


% ===== Insert after the limitations paragraph, before the Conclusion =========

A further extension concerns modality. We deliberately study tabular clinical
data, the most widely available modality across centers and the one on which
most deployed risk models are built, which also makes it the setting in which
the invariance assumptions can be stated and tested most cleanly. Oncology is
nonetheless moving toward decisions informed by imaging, digital pathology and
molecular profiling alongside the clinical record, and foundation models are
increasingly proposed to combine them <TODO-CITE: multimodal oncology FM>.
Multi-center shift is not removed by this move and may be sharpened by it, since
scanners, staining protocols and sequencing pipelines differ across sites in
ways that clinical variables do not. Whether a pretrained representation
transports the invariant signal that alignment fails to recover here is an open
question, and the diagnostic framing of this paper --- a known oracle, one
mechanism perturbed at a time --- transfers to that setting unchanged, since it
constrains the evaluation design rather than the model class.
